% Em. Lasker - J. R. Capablanca, St. Petersburg 1914: your comments. Format: games/README.md.
@ title | Em. Lasker - J. R. Capablanca, St. Petersburg 1914

@ 9...O-O | board
\mythoughts{White is in a must-win situation. These are the facts I have.

The pawn ending after all pieces are exchanged is winning for White, no matter who has the move:

\diagram{4k3/1pp2ppp/p1p5/8/4P3/8/PPP2PPP/4K3 w - - 0 1}

``If White can exchange all pieces, the pawn structure is a big advantage in the endgame.
Max Euwe gave the pure pawn ending (without pieces) resulting after the exchange of White's d-pawn
for Black's e-pawn as a win for White. The winning procedure is detailed in \emph{Secrets of Pawn Endings}.
In essence, the winning plan is to create a passed pawn on the kingside, while Black is unable to do
the same on the queenside because of the doubled pawns. The passed pawn ties down the black king and
allows the white king to transfer to the queenside at an opportune moment, forcing a pawn through to
promotion.'' (Wikipedia, ``Ruy Lopez, Exchange Variation'')
So exchanging pieces is probably good for White.

With the queens off the board, a mating attack is improbable. A possible winning plan is to create
a passed pawn on the kingside.

I do not know where to put my rooks. Doubling on the d-file only looks nice: the rooks have no way in,
because both d7 and d8 are covered.

Black has difficulty developing his light-squared bishop. ...f5 can be his next move.

I do not like 10.f4, which prepares my kingside passed pawn, because of 10...\B c5: my knight on d4
and my king on g1 stand on the same diagonal as the bishop.

I want to play a waiting move like 10.a4.}

@ 10...Re8
\mythoughts{I still do not like my king and knight on the same diagonal as the bishop.
I would play 11.\K h1.

The rook on e8 may prevent \B e3 tactically, since they are on the same file.}

@ 11...f6
\mythoughts{11...f6 is surprising for me: Black weakens the e6 square of his own free will.
12.f5 is logical, because it prepares a knight on e6. Black's bishop may keep an eye on that
square, but 12.f5 makes Black's development difficult.}

@ 12...b6
\mythoughts{I must develop my dark-squared bishop to connect my rooks.
Because of Black's rook on e8, I can never play the natural 13.\B e3: 13...\N xf5 14.exf5 \R xe3
and Black wins a pawn. 13.\B f4 is therefore logical.}
